% =====================================================================
%  ANÁLISIS DE REACCIONES: AMINAS FRENTE A AMIDAS
%  Se incluye desde Valentina/valentina.tex, antes de las conclusiones.
%  Valores calculados con Datos/datos_modelo_powerbi.xlsx (tabla Reacciones,
%  compuestos asignados); son los mismos KPI de la página
%  "6. Análisis de reacciones" del dashboard de Power BI.
% =====================================================================

\subsection{Análisis comparativo de reacciones: aminas frente a amidas}
\label{sec:analisis-reacciones}
La base de CAS-SciFinder no reporta rendimientos ni porcentajes de
conversión, así que la reactividad se comparó con las condiciones que cada
reacción necesitó: temperatura, tiempo y uso de catalizador
(Tabla~\ref{tab:reacciones}). Los porcentajes de temperatura se calculan
sobre las reacciones que reportan ese dato (61 de aminas y 38 de amidas).

\begin{table}[htbp]
	\centering
	\captionsetup{width=\columnwidth, justification=justified}
	\caption{Condiciones de reacción de las aminas y amidas asignadas.}
	\label{tab:reacciones}
	\small
	\begin{tabular}{lcc}
		\toprule
		\textbf{Indicador} & \textbf{Aminas} & \textbf{Amidas} \\
		\midrule
		Reacciones registradas                     & 65          & 40          \\
		Temperatura mediana (\si{\degreeCelsius})  & 25          & 95          \\
		Reacciones a $T \le 30$\,\si{\degreeCelsius} & \SI{68.9}{\percent} & \SI{23.7}{\percent} \\
		Reacciones a $T \ge 80$\,\si{\degreeCelsius} & \SI{11.5}{\percent} & \SI{57.9}{\percent} \\
		Tiempo mediano (h)                         & 3,5         & 12          \\
		Uso de catalizador                         & \SI{23.1}{\percent} & \SI{57.5}{\percent} \\
		Sustitución nucleofílica (SN)              & \SI{55.4}{\percent} & \SI{47.5}{\percent} \\
		\bottomrule
	\end{tabular}
\end{table}

Las aminas reaccionan sobre todo a temperatura ambiente y en pocas horas:
su par libre $sp^3$ es un buen nucleófilo y la sustitución nucleofílica es
el mecanismo dominante (Sección~\ref{sec:reac-aminas}). Las amidas, en
cambio, necesitan una temperatura mediana 70\,\si{\degreeCelsius} mayor,
más de tres veces el tiempo y más del doble de uso de catalizadores. Las
transamidaciones de la base (acetamida, N-metilacetamida y
N,N-dimetilacetamida con bencilamina) se hicieron entre 100 y
150\,\si{\degreeCelsius} con catalizadores de Fe, Cu, Nd o \ce{MoS2}, y las
condensaciones de la urea llegaron a 200\,\si{\degreeCelsius}.

Esta diferencia es la consecuencia práctica de la estabilidad
termodinámica descrita en la Sección~\ref{sec:estructura}: la resonancia
N--C=O hace que el carbonilo de la amida sea poco electrófilo y que su
nitrógeno sea poco nucleófilo, de modo que hay que suministrar energía (o
bajar la barrera con un catalizador) para romper o formar el enlace C--N.
La misma resonancia explica los momentos dipolares altos y los puntos de
ebullición elevados de las amidas en las Figuras~\ref{fig:spearman-int} y~\ref{fig:violin-cat}. Estos
indicadores pueden filtrarse por familia, medio y tipo de reacción en la
página ``Análisis de reacciones'' del dashboard de Power BI.
